% GENERATED FILE. Edit canonical lessons, then run npm run generate:books.
% canonical-source-hash: fnv1a64:5457afdf7a3b4402
% canonical-lessons: MW-C227-kuvo, MW-C227-udo, MW-C227-tisayo, MW-C227-rahat, MW-C227-dukal

\chapter{The Village Well}
\label{ch:mw-the-village-well}
\hlchaptermodality{227}

\begin{chapteropening}
\textbf{By the end of this chapter:} \emph{I can describe the well, a thirsty camel and a dry year.}

Complete the last lesson of chapter 227: I can describe the well, a thirsty camel and a dry year.
\end{chapteropening}

\section[\texorpdfstring{\mw{कूवो} (kūvo)}{कूवो (kūvo)}]{\mw{कूवो} (kūvo) — a well}

\label{lesson:MW-C227-kuvo}



\begin{quote}
Before the new one: say the Marwadi for dry land, then the Marwadi for laid waste.
\end{quote}

\subsection*{You'll want to know: \mw{कूवो}}
\textbf{\mw{कूवो}} — \emph{kūvo} — \textquotedblleft{}a well\textquotedblright{}.

\begin{quote}
\mw{गांव} \mw{मांय} \mw{एक} \mw{कूवो} \mw{है।}
\end{quote}

\begin{quote}
\emph{gā̃v mā̃y ek kūvo hai.}
\end{quote}

\begin{quote}
There's a well in the village.
\end{quote}

A masculine noun.

\begin{cousinweb}
\textbf{\mw{कूवो}} (\emph{kūvo}) is inherited from Sanskrit \emph{kūpa}, \textquotedblleft{}a hole, a pit, a well\textquotedblright{}; CDIAL lists Marwari \emph{kūvo}, a well without a masonry lining (CDIAL 3400).
\end{cousinweb}

\begin{grammarlens}[title={What you've built}]
One more word for this chapter.
\end{grammarlens}

\subsection*{Your turn}
\begin{itemize}
\raggedright
  \item \emph{Say it:} \emph{kūvo}
  \item \emph{Say it:} \emph{kūvo}, once more
  \item \emph{Say it:} \emph{gā̃v mā̃y ek kūvo hai.} — There's a well in the village.
  \item \emph{Recall:} say the Marwadi for dry land, then the Marwadi for laid waste, then say \emph{kūvo} again
\end{itemize}

\begin{tcolorbox}[breakable,colback=teal!4,colframe=teal!35!black,title={Before you move on}]
What does \mw{कूवो} mean? (\textquotedblleft{}A well\textquotedblright{}.) Say it once more.
\end{tcolorbox}
\section[\texorpdfstring{\mw{ऊंडो} (ū̃ḍo)}{ऊंडो (ū̃ḍo)}]{\mw{ऊंडो} (ū̃ḍo) — deep down; deep (of a well or a pit)}

\label{lesson:MW-C227-udo}



\begin{quote}
Before the new one: say the Marwadi for laid waste, then the Marwadi for a well.
\end{quote}

\subsection*{You'll want to know: \mw{ऊंडो}}
\textbf{\mw{ऊंडो}} — \emph{ū̃ḍo} — \textquotedblleft{}deep down; deep (of a well or a pit)\textquotedblright{}.

\begin{quote}
\mw{कूवो} \mw{घणो} \mw{ऊंडो} \mw{है।}
\end{quote}

\begin{quote}
\emph{kūvo ghaṇo ū̃ḍo hai.}
\end{quote}

\begin{quote}
The well is very deep.
\end{quote}

\begin{cousinweb}
is inherited from a reconstructed \emph{*ūṇḍa}, \textquotedblleft{}deep\textquotedblright{}; CDIAL lists Marwari \emph{ū̃ṛo} (CDIAL 2402). \textbf{\mw{गैरो}} is deep too; \textbf{\mw{ऊंडो}} is said especially of a well or a pit that goes far down.
\end{cousinweb}

\begin{grammarlens}[title={What you've built}]
One more word for this chapter.
\end{grammarlens}

\subsection*{Your turn}
\begin{itemize}
\raggedright
  \item \emph{Say it:} \emph{ū̃ḍo}
  \item \emph{Say it:} \emph{ū̃ḍo}, once more
  \item \emph{Say it:} \emph{kūvo ghaṇo ū̃ḍo hai.} — The well is very deep.
  \item \emph{Recall:} say the Marwadi for laid waste, then the Marwadi for a well, then say \emph{ū̃ḍo} again
\end{itemize}

\begin{tcolorbox}[breakable,colback=teal!4,colframe=teal!35!black,title={Before you move on}]
What does \mw{ऊंडो} mean? (\textquotedblleft{}Deep down; deep (of a well or a pit)\textquotedblright{}.) Say it once more.
\end{tcolorbox}
\section[\texorpdfstring{\mw{तिसायो} (tisāyo)}{तिसायो (tisāyo)}]{\mw{तिसायो} (tisāyo) — thirsty}

\label{lesson:MW-C227-tisayo}



\begin{quote}
Before the new one: say the Marwadi for a well, then the Marwadi for deep down.
\end{quote}

\subsection*{You'll want to know: \mw{तिसायो}}
\textbf{\mw{तिसायो}} — \emph{tisāyo} — \textquotedblleft{}thirsty\textquotedblright{}.

\begin{quote}
\mw{ऊंट} \mw{तिसायो} \mw{है।}
\end{quote}

\begin{quote}
\emph{ū̃ṭ tisāyo hai.}
\end{quote}

\begin{quote}
The camel is thirsty.
\end{quote}

\begin{cousinweb}
\textbf{\mw{तिसायो}} (\emph{tisāyo}) is inherited from Sanskrit \emph{tṛṣā}, \textquotedblleft{}thirst\textquotedblright{}; CDIAL lists Hindi \emph{tisāyā}, \textquotedblleft{}thirsty\textquotedblright{}, under it. CDIAL gives no Marwari form in this entry (CDIAL 5936).
\end{cousinweb}

\begin{grammarlens}[title={What you've built}]
One more word for this chapter.
\end{grammarlens}

\subsection*{Your turn}
\begin{itemize}
\raggedright
  \item \emph{Say it:} \emph{tisāyo}
  \item \emph{Say it:} \emph{tisāyo}, once more
  \item \emph{Say it:} \emph{ū̃ṭ tisāyo hai.} — The camel is thirsty.
  \item \emph{Recall:} say the Marwadi for a well, then the Marwadi for deep down, then say \emph{tisāyo} again
\end{itemize}

\begin{tcolorbox}[breakable,colback=teal!4,colframe=teal!35!black,title={Before you move on}]
What does \mw{तिसायो} mean? (\textquotedblleft{}Thirsty\textquotedblright{}.) Say it once more.
\end{tcolorbox}
\section[\texorpdfstring{\mw{रहट} (rahaṭ)}{रहट (rahaṭ)}]{\mw{रहट} (rahaṭ) — a Persian wheel (for raising water from a well)}

\label{lesson:MW-C227-rahat}



\begin{quote}
Before the new one: say the Marwadi for deep down, then the Marwadi for thirsty.
\end{quote}

\subsection*{You'll want to know: \mw{रहट}}
\textbf{\mw{रहट}} — \emph{rahaṭ} — \textquotedblleft{}a Persian wheel (for raising water from a well)\textquotedblright{}.

\begin{quote}
\mw{रहट} \mw{दिनूगै} \mw{चालै} \mw{है।}
\end{quote}

\begin{quote}
\emph{rahaṭ dinūgai chālai hai.}
\end{quote}

\begin{quote}
The Persian wheel turns early in the morning.
\end{quote}

A masculine noun.

\begin{cousinweb}
\textbf{\mw{रहट}} (\emph{rahaṭ}) is inherited from Sanskrit \emph{araghaṭṭa}, \textquotedblleft{}a wheel for raising water\textquotedblright{}; CDIAL lists Hindi \emph{rahaṭ}, \textquotedblleft{}a Persian wheel\textquotedblright{}, and a Gujarati word for a waterwheel. CDIAL gives no Marwari form in this entry (CDIAL 596).
\end{cousinweb}

\begin{grammarlens}[title={What you've built}]
One more word for this chapter.
\end{grammarlens}

\subsection*{Your turn}
\begin{itemize}
\raggedright
  \item \emph{Say it:} \emph{rahaṭ}
  \item \emph{Say it:} \emph{rahaṭ}, once more
  \item \emph{Say it:} \emph{rahaṭ dinūgai chālai hai.} — The Persian wheel turns early in the morning.
  \item \emph{Recall:} say the Marwadi for deep down, then the Marwadi for thirsty, then say \emph{rahaṭ} again
\end{itemize}

\begin{tcolorbox}[breakable,colback=teal!4,colframe=teal!35!black,title={Before you move on}]
What does \mw{रहट} mean? (\textquotedblleft{}A Persian wheel (for raising water from a well)\textquotedblright{}.) Say it once more.
\end{tcolorbox}
\section[\texorpdfstring{\mw{दुकाळ} (dukāḷ)}{दुकाळ (dukāḷ)}]{\mw{दुकाळ} (dukāḷ) — a drought; a famine}

\label{lesson:MW-C227-dukal}



\begin{quote}
Before the new one: say the Marwadi for thirsty, then the Marwadi for a Persian wheel.
\end{quote}

\subsection*{You'll want to know: \mw{दुकाळ}}
\textbf{\mw{दुकाळ}} — \emph{dukāḷ} — \textquotedblleft{}a drought; a famine\textquotedblright{}.

\begin{quote}
\mw{दुकाळ} \mw{मांय} \mw{पाणी} \mw{कोनी।}
\end{quote}

\begin{quote}
\emph{dukāḷ mā̃y pāṇī konī.}
\end{quote}

\begin{quote}
In a drought there's no water.
\end{quote}

A masculine noun.

\begin{cousinweb}
\textbf{\mw{दुकाळ}} (\emph{dukāḷ}) is inherited from Sanskrit \emph{duṣkāla}, \textquotedblleft{}an evil time\textquotedblright{}, built from \emph{duṣ-}, \textquotedblleft{}bad\textquotedblright{}, and \emph{kāla}, \textquotedblleft{}time\textquotedblright{}; CDIAL lists Old Marwari \emph{dūkāla}, \textquotedblleft{}famine\textquotedblright{} (CDIAL 6470).
\end{cousinweb}

\begin{grammarlens}[title={What you've built}]
One more word for this chapter.
\end{grammarlens}

\subsection*{Your turn}
\begin{itemize}
\raggedright
  \item \emph{Say it:} \emph{dukāḷ}
  \item \emph{Say it:} \emph{dukāḷ}, once more
  \item \emph{Say it:} \emph{dukāḷ mā̃y pāṇī konī.} — In a drought there's no water.
  \item \emph{Recall:} say the Marwadi for thirsty, then the Marwadi for a Persian wheel, then say \emph{dukāḷ} again
\end{itemize}

\begin{tcolorbox}[breakable,colback=teal!4,colframe=teal!35!black,title={Before you move on}]
What does \mw{दुकाळ} mean? (\textquotedblleft{}A drought; a famine\textquotedblright{}.) Say it once more.
\end{tcolorbox}
